\section{Choi positivity and exact Kraus-word spans}
\label{sec:kraus_iterate_choi}

Let $D\geq1$, let $(K_i)_{i\in I}$ be a finite family in $\MN{D}$, and define
\begin{align}
    E(X)&=\sum_{i\in I}K_iXK_i^\dagger.
    \notag
\end{align}
For a word $\sigma=(\sigma_0,\ldots,\sigma_{m-1})\in I^m$, write
\begin{align}
    K_\sigma&=K_{\sigma_0}\cdots K_{\sigma_{m-1}},
    &
    \mathcal S_m&=\operatorname{span}\{K_\sigma:\sigma\in I^m\}.
    \notag
\end{align}
The empty product is the identity matrix.

\begin{definition}[Eventually full vector spread]
    \label{def:kraus_eventually_full_vector_spread}
    \lean{Kraus.HasEventuallyFullVectorSpread}
    \leanok
    \uses{def:vector_spread_span}
    For $\psi\in\C^D$ and $m\in\mathbb N$, define
    \begin{align}
        H_m(K,\psi)
        &=\operatorname{span}\{K_\sigma\psi:\sigma\in I^m\}.
        \notag
    \end{align}
    The Kraus family has eventually full vector spread if, for every
    sufficiently large $m$ and every nonzero $\psi\in\C^D$,
    $H_m(K,\psi)=\C^D$. This is
    \cite[Theorem~6.8(2)]{Wolf2012Quantum}.
\end{definition}

\begin{definition}[Wolf's explicit eventual thresholds]
    \label{def:wolf_6_8_explicit_thresholds}
    \lean{Kraus.HasFullVectorSpreadFrom, Kraus.HasFullWordSpanFrom,
        Kraus.HasPosDefChoiFrom}
    \leanok
    \uses{def:kraus_eventually_full_vector_spread,
        def:kraus_eventually_full_word_span}
    For $n,q\in\mathbb N$, write:
    \begin{align}
        V_n(K)&:\Longleftrightarrow
          \forall m\geq n\ \forall\psi\neq0,\quad
          H_m(K,\psi)=\C^D, \notag\\
        W_q(K)&:\Longleftrightarrow
          \forall m\geq q,\quad \mathcal S_m=\MN{D}, \notag\\
        C_q(K)&:\Longleftrightarrow
          \forall m\geq q,\quad \tau_{E^m}>0.
        \notag
    \end{align}
    These definitions retain the universal quantifiers in items~2--4 of
    \cite[Theorem~6.8]{Wolf2012Quantum}; they are the explicit-threshold
    forms of the corresponding eventual predicates.
\end{definition}

\begin{theorem}[Full word span gives full vector spread]
    \label{thm:kraus_vector_spread_of_word_span}
    \lean{Kraus.vectorSpreadSpan_eq_top_of_wordSpan_eq_top}
    \leanok
    \uses{def:kraus_exact_word_span,def:vector_spread_span}
    If $\mathcal S_m=\MN{D}$, then
    $H_m(K,\psi)=\C^D$ for every nonzero $\psi\in\C^D$.
\end{theorem}

\begin{proof}
    \leanok
    Choose $k$ with $\psi_k\neq0$. For $v\in\C^D$, the matrix
    \begin{align}
        X&=\sum_{j=0}^{D-1}\frac{v_j}{\psi_k}E_{jk}
        \notag
    \end{align}
    satisfies $X\psi=v$. Hence the map $X\mapsto X\psi$ from $\MN{D}$ to
    $\C^D$ is surjective. If $\mathcal S_m=\MN{D}$, its image is precisely
    $H_m(K,\psi)$.
\end{proof}

\begin{theorem}[Eventually full word span gives eventual vector spread]
    \label{thm:kraus_eventually_full_vector_spread_of_word_span}
    \lean{Kraus.hasEventuallyFullVectorSpread_of_hasEventuallyFullWordSpan}
    \leanok
    \uses{def:kraus_eventually_full_word_span,
        def:kraus_eventually_full_vector_spread}
    If $\mathcal S_m=\MN{D}$ for every sufficiently large $m$, then the
    Kraus family has eventually full vector spread.
\end{theorem}

\begin{proof}
    \leanok
    \uses{thm:kraus_vector_spread_of_word_span}
    Apply Theorem~\ref{thm:kraus_vector_spread_of_word_span} at every
    sufficiently large length.
\end{proof}

\begin{theorem}[Positive-definite Choi matrix at a fixed iterate]
    \label{thm:kraus_iterate_choi_posdef_iff_word_span}
    \lean{Kraus.choiMatrix_mapLM_pow_posDef_iff_wordSpan_eq_top}
    \leanok
    \uses{thm:kraus_map_iterate_word_expansion,
        thm:choi_kraus_outer_product,
        thm:finite_frame_posdef}
    Let $\tau_{E^m}$ be the normalized Choi matrix of the $m$-fold iterate of
    $E$. Then
    \begin{align}
        \tau_{E^m}>0
        \quad\Longleftrightarrow\quad
        \mathcal S_m=\MN{D}.
        \label{eq:kraus_iterate_choi_posdef_iff_word_span}
    \end{align}
\end{theorem}

\begin{proof}
    \leanok
    \uses{thm:kraus_map_iterate_word_expansion,
        thm:choi_kraus_outer_product,
        thm:finite_frame_posdef}
    The Kraus operators of $E^m$ are the matrices $K_\sigma$ with
    $\sigma\in I^m$. Hence
    \begin{align}
        \tau_{E^m}
        &=\sum_{\sigma\in I^m}\ketbra{v_\sigma}{v_\sigma},
        &
        (v_\sigma)_{(a,b)}
        &=D^{-1/2}(K_\sigma)_{ab}.
        \notag
    \end{align}
    The matrix indices retain their row--column order. Since $D^{-1/2}\neq0$,
    the vectors $v_\sigma$ span $\C^{D\times D}$ exactly when the matrices
    $K_\sigma$ span $\MN{D}$. The finite-frame criterion now gives
    Equation~\eqref{eq:kraus_iterate_choi_posdef_iff_word_span}.
\end{proof}

\begin{corollary}[Eventual Choi positivity]
    \label{cor:kraus_iterate_eventual_choi_posdef}
    \lean{Kraus.eventually_choiMatrix_mapLM_pow_posDef_iff_hasEventuallyFullWordSpan}
    \leanok
    \uses{thm:kraus_iterate_choi_posdef_iff_word_span}
    The following conditions are equivalent:
    \begin{enumerate}
        \item $\tau_{E^m}>0$ for every sufficiently large $m$;
        \item $\mathcal S_m=\MN{D}$ for every sufficiently large $m$.
    \end{enumerate}
\end{corollary}

\begin{proof}
    \leanok
    \uses{thm:kraus_iterate_choi_posdef_iff_word_span}
    Apply
    Equation~\eqref{eq:kraus_iterate_choi_posdef_iff_word_span} pointwise for
    every sufficiently large $m$.
\end{proof}

\begin{theorem}[Primitive irreducible channels]
    \label{thm:kraus_primitive_irreducible_eventual_span_choi}
    \lean{Kraus.hasEventuallyFullWordSpan_of_isIrreducibleMap_of_isPrimitive,
        Kraus.eventually_choiMatrix_mapLM_pow_posDef_of_isIrreducibleMap_of_isPrimitive}
    \leanok
    \uses{def:kraus_eventually_full_word_span, def:irreducible_cp,
        def:primitive_channel}
    Suppose that $K$ is trace preserving and that its Kraus map $E$ is
    irreducible with peripheral spectrum $\{1\}$. This is the implication from
    item~1 to items~3 and~4 of
    \cite[Theorem~6.8(1,3,4)]{Wolf2012Quantum}. Then, for every sufficiently
    large $m$,
    \begin{align}
        \mathcal S_m&=\MN{D}.
        \notag
    \end{align}
    Moreover,
    \begin{align}
        \tau_{E^m}&>0.
        \notag
    \end{align}
\end{theorem}

\begin{proof}
    \leanok
    \uses{cor:kraus_iterate_eventual_choi_posdef}
    Irreducibility gives a positive-definite density matrix $\rho$ fixed by
    $E$. Put $N=E-P_\rho$, where
    $P_\rho(X)=\operatorname{tr}(X)\rho$, and define
    \begin{align}
        Q_F(B)&=\sum_{i,k}(B^\dagger F(e_{ik})B)_{ik}.
        \notag
    \end{align}
    For every word length $m$, expansion of the Kraus-map power gives
    \begin{align}
        \sum_{\sigma\in I^m}
        \left|\operatorname{tr}(B^\dagger K_\sigma)\right|^2
        &=\operatorname{Re}Q_{E^m}(B).
        \notag
    \end{align}
    The trivial peripheral spectrum implies $\lVert N^m\rVert\to0$, while
    $E^m=P_\rho+N^m$. Compactness of the unit sphere gives a constant
    $\delta>0$ such that
    \begin{align}
        \operatorname{Re}Q_{P_\rho}(B)&\geq
        \delta\lVert B\rVert^2,
        \notag
    \end{align}
    and the operator-norm estimate gives
    \begin{align}
        \left|\operatorname{Re}Q_{N^m}(B)\right|&\leq
        D^3\lVert N^m\rVert\lVert B\rVert^2.
        \notag
    \end{align}
    Hence $\operatorname{Re}Q_{E^m}(B)>0$ for every nonzero $B$ and all
    sufficiently large $m$. The trace-pairing identity then excludes a
    nonzero functional annihilating $\mathcal S_m$, so
    $\mathcal S_m=\MN{D}$. Finally,
    Corollary~\ref{cor:kraus_iterate_eventual_choi_posdef} gives the Choi
    statement.
\end{proof}

\begin{theorem}[Completely positive primitive maps]
    \label{thm:wolf_6_8_kraus_equivalences}
    \lean{Kraus.hasEventuallyFullVectorSpread_iff_exists_hasFullVectorSpreadFrom,
        Kraus.hasEventuallyFullWordSpan_iff_exists_hasFullWordSpanFrom,
        Kraus.eventually_choiMatrix_mapLM_pow_posDef_iff_exists_hasPosDefChoiFrom,
        Kraus.wolf_theorem_6_8_tfae,
        Kraus.wolf_theorem_6_8_minimal_indices}
    \leanok
    \uses{def:wolf_6_8_explicit_thresholds,
        thm:kraus_primitive_irreducible_eventual_span_choi,
        thm:kraus_iterate_choi_posdef_iff_word_span,
        thm:kraus_vector_spread_of_word_span,
        thm:kraus_irreducible_from_spreading,
        thm:kraus_primitive_from_tp_spreading}
    Let $K$ be trace preserving, so its Kraus map $E$ is completely positive
    and trace preserving. Then the following four statements are equivalent:
    \begin{enumerate}
        \item $E$ is irreducible and has trivial peripheral spectrum;
        \item $V_n(K)$ holds for some $n\in\mathbb N$;
        \item $W_q(K)$ holds for some $q\in\mathbb N$;
        \item $C_q(K)$ holds for some $q\in\mathbb N$.
    \end{enumerate}
    The first clause is exactly the source's term ``primitive'': Wolf's
    definition includes irreducibility, while the Lean predicate
    \texttt{IsPrimitive} records triviality of the peripheral spectrum.
    The admissible thresholds in clauses~3 and~4 coincide. Consequently one
    common minimal $q$ works in both clauses, and if $n$ is minimal in
    clause~2, then $n\leq q$. This is
    \cite[Theorem~6.8]{Wolf2012Quantum}.
\end{theorem}

\begin{proof}
    \leanok
    \uses{thm:kraus_primitive_irreducible_eventual_span_choi,
        thm:kraus_iterate_choi_posdef_iff_word_span,
        thm:kraus_vector_spread_of_word_span,
        thm:kraus_irreducible_from_spreading,
        thm:kraus_primitive_from_tp_spreading}
    Irreducibility and the trivial peripheral spectrum give $W_q(K)$ for a
    sufficiently large $q$. At each fixed length, vectorization identifies
    the span of the Kraus words with the span generating the iterate Choi
    matrix, so $W_q(K)\Longleftrightarrow C_q(K)$ with the same $q$.
    Moreover, a full matrix word span acts transitively on every nonzero
    vector, proving $W_q(K)\Rightarrow V_q(K)$. Conversely, evaluating
    $V_n(K)$ at $m=n$ gives both irreducibility and trivial peripheral
    spectrum by the fixed-length vector-spreading theorems. Thus the four
    clauses are equivalent. Pointwise equivalence makes the sets of
    admissible word and Choi thresholds identical. Since $q$ is itself an
    admissible vector threshold, minimality gives $n\leq q$.
\end{proof}

\begin{theorem}[Eventual Choi positivity gives eventual vector spread]
    \label{thm:kraus_eventually_full_vector_spread_of_choi_posdef}
    \lean{Kraus.hasEventuallyFullVectorSpread_of_eventually_choiMatrix_mapLM_pow_posDef}
    \leanok
    \uses{def:kraus_eventually_full_vector_spread}
    If $\tau_{E^m}>0$ for every sufficiently large $m$, then the Kraus family
    has eventually full vector spread.
\end{theorem}

\begin{proof}
    \leanok
    \uses{cor:kraus_iterate_eventual_choi_posdef,
        thm:kraus_eventually_full_vector_spread_of_word_span}
    Corollary~\ref{cor:kraus_iterate_eventual_choi_posdef} gives
    $\mathcal S_m=\MN{D}$ for every sufficiently large $m$. Apply
    Theorem~\ref{thm:kraus_eventually_full_vector_spread_of_word_span}.
\end{proof}
